% ===========================================================================
%  osc_experiments.tex
%
%  Numerical experiments on the oscillatory / high-order PINN suite, written
%  one-problem-at-a-time in the JCP/SIAM style of /root/paper/main.tex: each
%  family states its governing equation, exact solution and setting, and is then
%  immediately followed by ITS OWN result table, ITS OWN convergence figure
%  (relative L2 and loss vs training time), and a short discussion -- rather than
%  collecting all numbers in a single appendix-style dump.
%
%  Figures: docs/paper/figures/fig_<key>.pdf, generated by
%  experiments/plot_convergence.py from results/*_v3.csv (accuracy-vs-sweep
%  panel) and results/hist/*_history.json (convergence panels).
%
%  Class-agnostic: amsmath + booktabs + graphicx, \eqref/\ref, no cleveref.
% ===========================================================================

\section{Numerical experiments on oscillatory and high-order PDEs}\label{sec:osc}

Because \(\sinh\) is entire on \(\C\), the Taylor-jet backend of
Section~\ref{sec:impl} differentiates a \emph{complex-parameter} \(\sinh\)
network to arbitrary order with no dtype conversion in the inner loop. This makes
such a network usable as a PINN ansatz, and raises a concrete question:
\emph{at a matched real-parameter budget and a shared training protocol, does
carrying the network's value field in \(\C\) buy accuracy on oscillatory and
high-order PDEs?} We answer it on a literature-grounded suite of nine PDE
families, comparing the complex-parameter hyperbolic-sine network
(``complex \(\sinh\)'') against strong real high-frequency baselines. Each family
below states its problem, then reports its own result.

% ---------------------------------------------------------------------------
\subsection{Benchmark suite and shared protocol}\label{sec:bm:protocol}

Every problem is drawn from, or matched to, the PINN literature; the exact
manufactured solutions and parameters are stated with each family.
Table~\ref{tab:suite} summarizes the suite, which is organized into three groups:
\emph{high-frequency second-order} problems where spectral bias dominates
(Sections~\ref{sec:bm:helmholtz}--\ref{sec:bm:chirp}); \emph{high-order real}
operators where the derivative order dominates
(Sections~\ref{sec:bm:poly}--\ref{sec:bm:ch}); and \emph{genuinely
complex-valued} fields (Sections~\ref{sec:bm:nls}--\ref{sec:bm:maxwell}).

\begin{table}[htbp]
\centering
\small
\caption{Benchmark suite at a glance. Coordinates of a space--time problem are
written \((x,t)\); \(\Delta\) is the spatial Laplacian; \(a,k,m\) are the swept
parameters (wavenumber, mode, or temporal frequency). ``Source'' lists the
benchmark from which the setting is taken or to which it is matched.}
\label{tab:suite}
\begin{tabular}{@{}l l c l l@{}}
\toprule
family & operator & order & exact solution & source \\
\midrule
Helmholtz (high \(k\)) & \(\Delta u + \kappa^2 u\) & 2 & \(\sin(a\pi x)\sin(a\pi y)\) & \cite{Wang2021Pathologies} \\
Helmholtz (variable \(\kappa\)) & \(\Delta u + \kappa^2(x)u\) & 2 & \(\sin(a\pi x)\sin(a\pi y)\) & \cite{Hao2024PINNacle} \\
radial chirp & \(-\Delta u + u\) & 2 & \(\sin\!\big(\tfrac{a\pi}{2}r^2\big)\) & \cite{Tancik2020Fourier,Liu2020MscaleDNN} \\
polyharmonic & \(\Delta^m u\) & \(2m\) & \(\sin(\pi x)\,[\sin(\pi y)]\) & \cite{Vahab2022Biharmonic} \\
plate / beam & \(\Delta^2 w\) & 4 & \(\sin(m\pi x)\sin(n\pi y)\) & \cite{Vahab2022Biharmonic} \\
linearized KdV & \(u_t + \delta\,u_{xxx}\) & 3 & \(\sin(k\pi x)\cos(k\pi t)\) & \cite{Raissi2019} \\
Cahn--Hilliard & \(u_t - M\Delta(u^3-u-\gamma\Delta u\,[+\kappa\Delta^2 u])\) & 4/6 & \(\sin(a\pi x)\cos(a\pi t)\) & \cite{Raissi2019,Hao2024PINNacle} \\
cubic NLS & \(\mathrm{i}u_t + \tfrac12 u_{xx} + |u|^2 u\) & 2 & \(\operatorname{sech}(x)\,\mathrm{e}^{\mathrm{i}kt}\) & \cite{Raissi2019} \\
Maxwell (lossy) & \(\Delta E + \kappa^2 E\), \(\kappa^2\in\mathbb{C}\) & 2 & \(\mathrm{e}^{\mathrm{i}a\pi(x+y)}\) & \cite{Jiang2024Maxwell} \\
\bottomrule
\end{tabular}
\end{table}

\paragraph{Architectures.}
The complex-sinh network is a multilayer perceptron with \texttt{complex128}
parameters and the entire activation \(\sinh\). It is compared against four real
``high-frequency'' baselines that are standard remedies for spectral bias:
a random Fourier-feature network~\cite{Tancik2020Fourier}, a sinusoidal
representation network (SIREN)~\cite{Sitzmann2020SIREN}, a multi-scale deep
network (MscaleDNN)~\cite{Liu2020MscaleDNN}, and plain real \(\tanh\) and
\(\sinh\) baselines. For genuinely complex-valued targets
(Sections~\ref{sec:bm:nls}--\ref{sec:bm:maxwell}) each real baseline carries the
field as a split real--imaginary pair (an \(\mathbb{R}\)-valued PINN, RVPINN),
as in the original complex-valued PINN formulation~\cite{Raissi2019}.

\paragraph{Parameter and frequency matching.}
A width-\(H\) complex network has \(\approx 2H^2\) real degrees of freedom per
layer, so every single real baseline is widened to
\(H_{\mathbb{R}} = \lceil \sqrt{2}\,H \rceil\) (giving it at least the complex
network's real parameter count); each component of a split-real pair uses width
\(H\) so that the \emph{pair} matches. Frequency-aware architectures are
initialized at the problem's characteristic wavenumber -- the SIREN /
complex-sinh first-layer scale \(\omega_0\) and the Fourier bandwidth \(\sigma\)
are set proportional to the target spatial frequency of each case -- so each
method receives its best start at every swept frequency. For the high-order
families \(\omega_0\) is held \emph{at} the target frequency rather than above
it, because an order-\(m\) operator amplifies an initialization frequency
\(\omega\) like \(\omega^{m}\), so an over-large \(\omega_0\) buries the learning
signal.

\paragraph{Derivative backend and nonlinear fluxes.}
All jet-compatible architectures evaluate every residual derivative with the
\emph{same} complex-Waring Taylor-jet operator of Section~\ref{sec:impl}, so the
derivative machinery does not bias the comparison. Nonlinear fluxes are written
purely from single-monomial partials of the network: e.g.\ the Cahn--Hilliard
term
\begin{equation}\label{eq:lap_u3}
  \Delta(u^3) \;=\; 3u^2\,u_{xx} + 6u\,(u_x)^2
\end{equation}
uses only \(u\), \(u_x\), \(u_{xx}\), each a single Taylor-jet evaluation.

\paragraph{Optimization and metric.}
Training uses Adam with a wall-clock budget shared across all methods, a
time-based cosine learning-rate schedule decaying from \(10^{-3}\) to a floor of
\(10^{-4}\), and a fixed dense collocation set. Boundary and initial data are
imposed by a penalty term of weight \(100\); for the polyharmonic and plate
families the simply-supported (Navier) conditions are enforced on the lower even
derivatives. The reported accuracy is the relative \(L^2\) error of
\(\operatorname{Re} u\) (or of \(u\) for complex-valued targets) against the
exact solution on a held-out grid of \(8192\) points,
\begin{equation}\label{eq:relL2}
  \mathcal{E}_{L^2}
  \;=\;
  \frac{\big\|\,u_\theta - u_\star\,\big\|_{L^2(\Omega)}}
       {\big\|\,u_\star\,\big\|_{L^2(\Omega)}} ,
\end{equation}
averaged over independent random seeds. For each family we also report the
\emph{advantage factor}
\begin{equation}\label{eq:adv}
  \mathrm{adv}
  \;=\;
  \frac{\min_{\text{real baselines}} \mathcal{E}_{L^2}}
       {\mathcal{E}_{L^2}(\text{complex }\sinh)} ,
\end{equation}
so \(\mathrm{adv}>1\) means the complex network beats the best real contender on
that instance.

\paragraph{Reading the figures.}
Each family is accompanied by a three-panel figure: panel~(a) plots the relative
\(L^2\) error against the swept parameter for all methods (the same numbers as
the family's table); panels~(b) and~(c) plot, respectively, the relative \(L^2\)
error and the interior residual loss against \emph{training time} for a
representative case, with shaded bands showing the seed-to-seed min--max range.
A dotted line at \(\mathcal{E}_{L^2}=1\) marks the trivial-predictor level;
curves at or above it indicate failure or divergence within the budget.

\paragraph{Hardware and budget.}
All runs are measured on a single NVIDIA H20 GPU (\(97\)~GB, CUDA 12.1,
PyTorch 2.5) in \texttt{float64}/\texttt{complex128}. Each (problem, variant,
seed) receives the same wall-clock budget (\(70\)~s) and the cosine schedule
above; reported errors are means over two seeds. Raw per-seed CSV files and the
convergence-trace sidecars are in \path{results/} and \path{results/hist/}.

% ===========================================================================
\subsection{High-wavenumber Helmholtz}\label{sec:bm:helmholtz}

The first benchmark is the constant-coefficient Helmholtz equation on the unit
square,
\begin{equation}\label{eq:helm_pde}
  \Delta u + \kappa^2 u = f
  \qquad \text{in } \Omega = (-1,1)^2,
  \qquad \kappa = a\pi,
\end{equation}
with homogeneous Dirichlet data \(u = 0\) on \(\partial\Omega\). The exact
solution is the separable eigenmode
\begin{equation}\label{eq:helm_sol}
  u_\star(x,y) = \sin(a\pi x)\,\sin(a\pi y),
\end{equation}
and the source follows in closed form,
\begin{align}\label{eq:helm_source}
  f
  &= \Delta u_\star + \kappa^2 u_\star
   = \big(-2(a\pi)^2 + (a\pi)^2\big)\,u_\star
   = -(a\pi)^2\,u_\star .
\end{align}
The wavenumber is swept over \(a \in \{2,4,6,8,10\}\), so the number of
half-periods across the domain grows while the operator order stays fixed at
two. This is the Helmholtz benchmark of Wang, Teng and
Perdikaris~\cite{Wang2021Pathologies}; the same family is the
Poisson--Boltzmann/Helmholtz task in PINNacle~\cite{Hao2024PINNacle}. We use the
isotropic case \(a_1=a_2=a\), \(\kappa=a\pi\), and additionally report the
anisotropic case \((a_1,a_2)=(1,4)\) of~\cite{Wang2021Pathologies} verbatim
(\(\kappa=\pi\), \(q=(\kappa^2-(a_1\pi)^2-(a_2\pi)^2)u_\star\)).
\emph{Setting:} \(\Omega=(-1,1)^2\); Dirichlet \(u=0\);
\(\omega_0=\max(10,2\pi a)\), \(\sigma=\max(2,\pi a)\); residual scale
\((a\pi)^2\).

\begin{table}[htbp]
\centering
\small
\caption{High-wavenumber Helmholtz: relative \(L^2\) error \eqref{eq:relL2}
(best in bold). The bottom row is the anisotropic \((1,4)\) case.}
\label{tab:res-helmholtz}
\begin{tabular}{@{}lrrrr@{}}
\toprule
\(a\) & complex \(\sinh\) & Fourier & SIREN & MscaleDNN \\
\midrule
2  & \textbf{5.22e-2} & 1.52e-1 & 8.09e-2 & 6.99e-1 \\
4  & \textbf{3.22e-1} & 4.32e-1 & 3.46e-1 & 1.00e0 \\
6  & \textbf{5.71e-1} & 8.51e-1 & 8.19e-1 & 1.00e0 \\
8  & \textbf{6.29e-1} & 9.32e-1 & 9.27e-1 & 1.00e0 \\
10 & \textbf{7.71e-1} & 9.86e-1 & 9.83e-1 & 1.00e0 \\
\midrule
\((1,4)\) & 3.89e-1 & 7.79e-1 & \textbf{2.21e-1} & 1.00e0 \\
\bottomrule
\end{tabular}
\end{table}

\begin{figure}[htbp]
  \centering
  \includegraphics[width=\linewidth]{fig_helmholtz}
  \caption{High-wavenumber Helmholtz. (a) Relative \(L^2\) error versus
  wavenumber \(a\); (b)--(c) relative \(L^2\) error and interior residual loss
  versus training time at \(a=6\). Bands show the seed min--max.}
  \label{fig:helmholtz}
\end{figure}

The complex network is the most accurate method at every isotropic wavenumber,
and its error grows the most slowly as \(a\) increases: at \(a=10\) it reaches
\(0.77\) while both SIREN and Fourier have saturated near \(0.98\) and MscaleDNN
has collapsed to the trivial level. The convergence traces (Figure~\ref{fig:helmholtz}b)
show the same ordering in time -- the complex curve separates early and keeps
descending. The single anisotropic \((1,4)\) case is the exception: at this low,
strongly directional frequency a SIREN reaches \(2.2\times10^{-1}\) against the
complex network's \(3.9\times10^{-1}\), consistent with SIREN's strength on
separable low-frequency sinusoids.
\FloatBarrier

% ===========================================================================
\subsection{Variable-coefficient (scattering) Helmholtz}\label{sec:bm:helmvc}

Moving from a single clean eigenmode to a heterogeneous medium, we let the
zeroth-order coefficient vary in space,
\begin{equation}\label{eq:helmvc_pde}
  \Delta u + \kappa^2(x)\,u = f
  \qquad \text{in } \Omega = (-1,1)^2,
\end{equation}
with a smooth lens-like contrast
\begin{equation}\label{eq:helmvc_kappa}
  \kappa^2(x,y) = (a\pi)^2\big(1 + C\,\sin(\pi x)\sin(\pi y)\big),
  \qquad C = 0.5 .
\end{equation}
The manufactured solution and Dirichlet data are again \eqref{eq:helm_sol} with
\(u=0\) on \(\partial\Omega\), and the source is
\(f = -2(a\pi)^2\,u_\star + \kappa^2(x)\,u_\star\). The background wavenumber is
swept over \(a \in \{2,4,6\}\). Variable-coefficient (scattering) Helmholtz in a
heterogeneous medium is the realistic electromagnetics/acoustics regime of the
PINNacle Poisson--Boltzmann/Helmholtz tasks~\cite{Hao2024PINNacle}; the
\(\pm 50\%\) smooth contrast here is a controlled, exactly-known-solution
instance that also exercises the variable-coefficient residual path
(\(\kappa^2(x)\) precomputed once on the collocation set).
\emph{Setting:} \(\Omega=(-1,1)^2\); Dirichlet \(u=0\); \(C=0.5\);
\(\omega_0,\sigma\) as in Section~\ref{sec:bm:helmholtz}.

\begin{table}[htbp]
\centering
\small
\caption{Variable-coefficient Helmholtz: relative \(L^2\) error \eqref{eq:relL2}.}
\label{tab:res-helmvc}
\begin{tabular}{@{}lrrrr@{}}
\toprule
\(a\) & complex \(\sinh\) & Fourier & SIREN & MscaleDNN \\
\midrule
2 & 3.59e-1 & \textbf{2.11e-1} & 3.15e-1 & 1.00e0 \\
4 & 4.64e-1 & 5.45e-1 & \textbf{4.37e-1} & 1.00e0 \\
6 & \textbf{7.12e-1} & 8.72e-1 & 8.07e-1 & 1.00e0 \\
\bottomrule
\end{tabular}
\end{table}

\begin{figure}[htbp]
  \centering
  \includegraphics[width=\linewidth]{fig_helmvc}
  \caption{Variable-coefficient Helmholtz. (a) Relative \(L^2\) versus
  wavenumber \(a\); (b)--(c) relative \(L^2\) and loss versus training time at
  \(a=6\). Bands show the seed min--max.}
  \label{fig:helmvc}
\end{figure}

Heterogeneity narrows the gap. At the lowest background wavenumber a
Fourier-feature network (\(a=2\)) and a SIREN (\(a=4\)) are ahead, and the
complex network leads only at the highest wavenumber \(a=6\), where its advantage
returns to \(1.13\times\). The trend is monotone in the right direction --
the advantage rises from \(0.59\) to \(1.13\) as \(a\) grows -- so the complex
network's relative benefit is again frequency-driven, but at low frequency the
smooth coefficient does not yet stress the real baselines enough to separate
them.
\FloatBarrier

% ===========================================================================
\subsection{Non-separable radial chirp}\label{sec:bm:chirp}

The eigenmode benchmarks are separable products of sines, which sine-based real
baselines can represent almost natively and which may therefore mask any
expressivity gap. To remove this confound we use a non-separable target with a
space-varying local frequency,
\begin{equation}\label{eq:chirp_pde}
  -\Delta u + u = f \qquad \text{in } \Omega = (-1,1)^2,
  \qquad
  u_\star(x,y) = \sin\!\big(\tfrac{a\pi}{2}(x^2+y^2)\big),
\end{equation}
a radial chirp whose instantaneous frequency \(|\nabla\phi| = a\pi r\) grows with
radius \(r=\sqrt{x^2+y^2}\). The source is assembled from
\begin{align}\label{eq:chirp_source}
  \Delta u_\star
  &= -(a\pi)^2 r^2 \sin\phi + 2a\pi\cos\phi, \quad
  f = -\Delta u_\star + u_\star, \quad \phi = \tfrac{a\pi}{2}r^2,
\end{align}
with Dirichlet data \(u=u_\star\) and sweep \(a\in\{2,4,6,8\}\). Space-varying
frequency targets are the standard stress test for high-frequency representation
methods such as Fourier features~\cite{Tancik2020Fourier} and multi-scale
networks~\cite{Liu2020MscaleDNN}, which are tuned for stationary spectra; this
benchmark is the fairness control of the suite, since every architecture must
\emph{construct} the varying oscillation rather than read it off a fixed
frequency basis.
\emph{Setting:} \(\Omega=(-1,1)^2\); Dirichlet \(u=u_\star\);
\(\omega_0,\sigma\) as in Section~\ref{sec:bm:helmholtz}.

\begin{table}[htbp]
\centering
\small
\caption{Radial chirp: relative \(L^2\) error \eqref{eq:relL2}. Errors above
one indicate divergence within the budget.}
\label{tab:res-chirp}
\begin{tabular}{@{}lrrrr@{}}
\toprule
\(a\) & complex \(\sinh\) & Fourier & SIREN & MscaleDNN \\
\midrule
2 & \textbf{3.76e-1} & 7.07e-1 & 8.74e-1 & 7.93e-1 \\
4 & \textbf{4.52e-1} & 9.19e-1 & 7.16e-1 & 1.60e0 \\
6 & \textbf{8.47e-1} & 1.18e0  & 9.76e-1 & 1.89e0 \\
8 & \textbf{8.30e-1} & 1.24e0  & 9.64e-1 & 2.39e0 \\
\bottomrule
\end{tabular}
\end{table}

\begin{figure}[htbp]
  \centering
  \includegraphics[width=\linewidth]{fig_chirp}
  \caption{Non-separable radial chirp. (a) Relative \(L^2\) versus chirp rate
  \(a\); (b)--(c) relative \(L^2\) and loss versus training time at \(a=4\).
  Bands show the seed min--max.}
  \label{fig:chirp}
\end{figure}

This is the cleanest expressivity signal in the suite. Because the solution is
not a product of one-dimensional sines, the sine-based MscaleDNN -- competitive
on separable targets -- degrades sharply, with errors above one (i.e.\ it
diverges), and the Fourier-feature network does the same from \(a=6\) on. The
complex network retains the smallest error throughout and is the only method that
stays below one at every chirp rate, confirming that its advantage is one of
representation, not of an aligned fixed basis.
\FloatBarrier

% ===========================================================================
\subsection{Polyharmonic order sweep}\label{sec:bm:poly}

To isolate the effect of the differential \emph{order} from that of frequency,
we fix the solution and frequency and vary only the order of a polyharmonic
operator. The eigenmode identity \(\Delta^m u = (-S)^m u\) holds for the
sinusoidal eigenfunctions below, giving
\begin{equation}\label{eq:poly_pde}
  \Delta^m u = (-S)^m\,u,
  \qquad
  \begin{cases}
    u_\star = \sin(\pi x)\sin(\pi y), & S = 2\pi^2 \ (\text{2D}),\\[2pt]
    u_\star = \sin(\pi x),            & S = \pi^2 \ (\text{1D}),
  \end{cases}
\end{equation}
with simply-supported (Navier) boundary conditions imposed exactly through the
lower even derivatives, \(\Delta^j u = 0\) on \(\partial\Omega\),
\(j=0,\dots,m-1\). The order is swept over \(2m\in\{2,4,6\}\) in two dimensions
and \(2m\in\{2,4,6,8,10\}\) in one dimension: in 1D the operator is a single
monomial \(\partial_x^{2m}\), cheap enough to reach order ten, whereas in 2D
\(\Delta^m\) expands into \(m+1\) high-order terms. Biharmonic and polyharmonic
operators are the canonical high-order PINN testbed~\cite{Vahab2022Biharmonic};
holding \(u_\star\), the frequency and the domain fixed while varying only \(m\)
is our isolation design for the order axis.
\emph{Setting:} \(\Omega=(-1,1)^d\), \(d\in\{1,2\}\); Navier BC; 1D
\(\omega_0=\sigma=\pi\), 2D \(\omega_0=2\pi,\sigma=\pi\); residual scale \(S^m\).

\begin{table}[htbp]
\centering
\small
\caption{Polyharmonic order sweep: relative \(L^2\) error \eqref{eq:relL2}.
The plain real \(\sinh\) baseline (omitted) exceeds \(0.85\) on every row with
\(2m\ge4\).}
\label{tab:res-poly}
\begin{tabular}{@{}llrrrr@{}}
\toprule
dim & \(2m\) & complex \(\sinh\) & Fourier & SIREN & MscaleDNN \\
\midrule
1D & 2  & \textbf{4.15e-5} & 2.33e-4 & 2.13e-4 & 3.09e-4 \\
   & 4  & \textbf{2.03e-3} & 1.20e-2 & 3.15e-2 & 6.83e-2 \\
   & 6  & 1.42e-1 & 6.90e-1 & 9.66e-1 & \textbf{1.25e-1} \\
   & 8  & \textbf{2.14e-1} & 1.00e0  & 1.00e0  & 4.08e-1 \\
   & 10 & \textbf{9.95e-1} & 1.00e0  & 1.01e0  & 9.98e-1 \\
\midrule
2D & 2  & \textbf{4.23e-4} & 6.66e-4 & 7.83e-4 & 1.09e-2 \\
   & 4  & 3.96e-2 & \textbf{3.17e-2} & 9.83e-1 & 8.99e-1 \\
   & 6  & \textbf{7.67e-2} & 9.99e-1 & 9.94e-1 & 9.96e-1 \\
\bottomrule
\end{tabular}
\end{table}

\begin{figure}[htbp]
  \centering
  \includegraphics[width=\linewidth]{fig_poly1d}\\[2pt]
  \includegraphics[width=\linewidth]{fig_poly2d}
  \caption{Polyharmonic order sweep, 1D (top) and 2D (bottom). (a) Relative
  \(L^2\) versus operator order \(2m\); (b)--(c) relative \(L^2\) and loss
  versus training time at the representative order (\(2m=4\) in 1D, \(2m=6\) in
  2D). Bands show the seed min--max.}
  \label{fig:poly}
\end{figure}

This is the family the ansatz is designed for, and the order axis bears it out.
On the single-monomial 1D operator the complex network is one to two orders of
magnitude more accurate than every real baseline at orders two and four, holds
the lead at order eight, and is edged out by MscaleDNN only at order six; at
order ten the fixed budget defeats every method. In two dimensions the gap
\emph{widens} with order, reaching a \(13\times\) advantage at the order-six
triharmonic where all real baselines have effectively failed
(Figure~\ref{fig:poly}a, bottom). The convergence traces show that the real
baselines plateau almost immediately at high order while the complex network
keeps reducing the residual, which is the behaviour predicted by the
rank-optimal complex schedule on balanced repeated-index operators.
\FloatBarrier

% ===========================================================================
\subsection{Kirchhoff plate and Euler--Bernoulli beam}\label{sec:bm:plate}

The fourth-order vibration eigenproblems of structural mechanics give three
oscillatory biharmonic benchmarks. For the Kirchhoff plate,
\begin{equation}\label{eq:plate_pde}
  \Delta^2 w = S^2\,w \quad \text{in } (-1,1)^2,
  \quad S = (m^2+n^2)\pi^2,
  \quad w_\star = \sin(m\pi x)\sin(n\pi y),
\end{equation}
with simply-supported edges \(w=0\), \(\Delta w=0\); for the Euler--Bernoulli
beam,
\begin{equation}\label{eq:beam_pde}
  w'''' = (m\pi)^4\,w \quad \text{in } (-1,1),
  \quad w_\star = \sin(m\pi x),
\end{equation}
with \(w=0\), \(w''=0\). We use the isotropic plate modes \((m,m)\),
\(m\in\{1,2,3\}\), the anisotropic mixed modes \((m,m{+}1)\), \(m\in\{2,3,4\}\),
and the beam modes \(m\in\{1,2,3\}\); the order is fixed at four and the sweep
raises only the oscillation. The simply-supported plate
(\(D\Delta^2 w=q\)) and beam are standard PINN structural
benchmarks~\cite{Vahab2022Biharmonic}; we use the free-vibration eigenmode form
with the matching manufactured right-hand side.
\emph{Setting:} Navier BC throughout; \(\omega_0=\max(10,2\pi f)\),
\(\sigma=\max(2,\pi f)\), \(f=\max(m,n)\); residual scale \(S^2\).

\begin{table}[htbp]
\centering
\small
\caption{Plate and beam (4th order): relative \(L^2\) error \eqref{eq:relL2}.
``mixed'' is the anisotropic plate mode \((m,m{+}1)\).}
\label{tab:res-plate}
\begin{tabular}{@{}llrrrr@{}}
\toprule
family & \(m\) & complex \(\sinh\) & Fourier & SIREN & MscaleDNN \\
\midrule
plate (iso) & 1 & 4.62e-2 & \textbf{4.24e-2} & 9.88e-1 & 9.80e-1 \\
            & 2 & \textbf{2.47e-1} & 2.76e-1 & 5.35e-1 & 1.00e0 \\
            & 3 & \textbf{3.63e-1} & 6.72e-1 & 9.24e-1 & 1.00e0 \\
\midrule
plate (mixed) & 2 & \textbf{6.59e-1} & 9.78e-1 & 9.39e-1 & 1.00e0 \\
              & 3 & \textbf{5.54e-1} & 9.93e-1 & 9.81e-1 & 1.00e0 \\
              & 4 & \textbf{6.90e-1} & 9.97e-1 & 9.88e-1 & 1.00e0 \\
\midrule
beam (1D) & 1 & 4.36e-2 & \textbf{1.35e-2} & 1.01e0  & 9.78e-2 \\
          & 2 & \textbf{1.29e-1} & 1.57e-1 & 6.11e-1 & 5.49e-1 \\
          & 3 & 2.22e-1 & \textbf{1.16e-1} & 8.06e-1 & 1.68e0 \\
\bottomrule
\end{tabular}
\end{table}

\begin{figure}[htbp]
  \centering
  \includegraphics[width=0.49\linewidth]{fig_plate}
  \hfill
  \includegraphics[width=0.49\linewidth]{fig_beam}\\[2pt]
  \includegraphics[width=\linewidth]{fig_platemix}
  \caption{Fourth-order plate/beam. Top: isotropic plate (left) and beam
  (right); bottom: anisotropic mixed-mode plate. In each, (a) relative \(L^2\)
  versus mode \(m\) and (b)--(c) the time traces at \(m=2\) (plate, beam) or
  modes \((3,4)\) (mixed). Bands show the seed min--max.}
  \label{fig:plate}
\end{figure}

The anisotropic (mixed-mode) plate is the decisive case: with two different
per-axis frequencies at fixed order four it is non-separable in the same sense as
the chirp, and the complex network wins every mode by a clear margin
(\(1.4\)--\(1.8\times\)) while all real baselines sit near the trivial level.
The isotropic plate is won at the higher modes, with an advantage growing from
parity at \(m=1\) to \(1.85\times\) at \(m=3\). The one-dimensional beam is the
subtler case: a Fourier-feature network ties or wins on the smoothest modes,
because a 1D fourth-order sine is precisely the stationary single-frequency
signal Fourier features are built for; the complex network still wins the
intermediate mode \(m=2\).
\FloatBarrier

% ===========================================================================
\subsection{Linearized Korteweg--de Vries / dispersive wave}\label{sec:bm:kdv}

The third-order dispersive operator is exercised by a linearized KdV (Airy)
equation in space--time,
\begin{equation}\label{eq:kdv_pde}
  u_t + \delta\,u_{xxx} = f \quad \text{in } (x,t)\in(-1,1)^2,
  \quad \delta = 1,
  \quad u_\star = \sin(k\pi x)\cos(k\pi t),
\end{equation}
with Dirichlet data \(u=u_\star\) on the space--time boundary and source
\begin{align}\label{eq:kdv_source}
  f
  &= u_{\star,t} + \delta\,u_{\star,xxx}
   = -k\pi\sin(k\pi x)\sin(k\pi t) - (k\pi)^3\cos(k\pi x)\cos(k\pi t),
\end{align}
swept over \(k\in\{2,3,4,5,6\}\). KdV is a canonical PINN
benchmark~\cite{Raissi2019}; we use the linear dispersive (Airy) operator with a
manufactured oscillatory solution to isolate the odd third-order term
\(u_{xxx}\) from the convective nonlinearity.
\emph{Setting:} \((x,t)\in(-1,1)^2\); Dirichlet \(u=u_\star\); \(\delta=1\);
\(\omega_0=\max(10,2\pi k)\), \(\sigma=\max(2,\pi k)\).

\begin{table}[htbp]
\centering
\small
\caption{Linearized KdV (3rd order): relative \(L^2\) error \eqref{eq:relL2}.
This is a counter-case: SIREN leads at the higher wavenumbers.}
\label{tab:res-kdv}
\begin{tabular}{@{}lrrrr@{}}
\toprule
\(k\) & complex \(\sinh\) & Fourier & SIREN & MscaleDNN \\
\midrule
2 & 3.65e-1 & \textbf{2.89e-1} & 4.54e-1 & 1.29e0 \\
3 & 5.68e-1 & 6.39e-1 & \textbf{6.05e-1} & 2.62e0 \\
4 & 8.70e-1 & 7.22e-1 & \textbf{6.55e-1} & 2.51e0 \\
5 & 9.48e-1 & 1.06e0  & \textbf{6.82e-1} & 4.03e0 \\
6 & 8.03e-1 & 1.10e0  & \textbf{7.37e-1} & 3.12e0 \\
\bottomrule
\end{tabular}
\end{table}

\begin{figure}[htbp]
  \centering
  \includegraphics[width=\linewidth]{fig_kdv}
  \caption{Linearized KdV. (a) Relative \(L^2\) versus wavenumber \(k\);
  (b)--(c) relative \(L^2\) and loss versus training time at \(k=4\). Bands show
  the seed min--max.}
  \label{fig:kdv}
\end{figure}

We report this family as an honest counter-case. The solution is a separable
pure sine of only third order -- exactly the regime of SIREN -- and SIREN is
indeed the most accurate at four of the five wavenumbers, the complex network
leading only marginally at \(k=3\); MscaleDNN diverges throughout. A
first-order-in-time dispersive operator does not present the balanced
repeated-index structure that the complex ansatz exploits, and the separable
sine removes the expressivity gap that the chirp exposes. This is precisely the
confound the chirp benchmark (Section~\ref{sec:bm:chirp}) is designed to control
for.
\FloatBarrier

% ===========================================================================
\subsection{Cahn--Hilliard}\label{sec:bm:ch}

The phase-field Cahn--Hilliard equation combines a high-order linear operator
with a cubic nonlinearity. In one space dimension plus time,
\begin{align}\label{eq:ch_pde}
  u_t &= M\big[\Delta(u^3) - \Delta u - \gamma\,\Delta^2 u\big] && (\text{4th order}),\\
  u_t &= M\big[\Delta(u^3) - \Delta u - \gamma\,\Delta^2 u + \kappa\,\Delta^3 u\big] && (\text{6th order}),
  \label{eq:ch_pde6}
\end{align}
with \(M=\gamma=\kappa=1\), \(\Delta=\partial_x^2\), \((x,t)\in(-1,1)^2\), and
manufactured solution \(u_\star = \sin(a\pi x)\cos(a\pi t)\). The nonlinear flux
is expanded as in \eqref{eq:lap_u3}; the source and the boundary targets (on
\(u\), \(u_{xx}\), and additionally \(u_{xxxx}\) at sixth order) come from the
analytic solution, with \(a\in\{2,3\}\) at orders \(4\) and \(6\).
Cahn--Hilliard is a standard fourth-order phase-field PINN benchmark and a
sixth-order variant appears in randomized-smoothing PINNs; such nonlinear
high-order space--time PDEs form the hardest tier of
PINNacle~\cite{Hao2024PINNacle}.
\emph{Setting:} \((x,t)\in(-1,1)^2\); \(M=\gamma=\kappa=1\);
\(\omega_0=\max(10,2\pi a)\), \(\sigma=\max(2,\pi a)\).

\begin{table}[htbp]
\centering
\small
\caption{Cahn--Hilliard (4th/6th order, nonlinear): relative \(L^2\) error
\eqref{eq:relL2}.}
\label{tab:res-ch}
\begin{tabular}{@{}llrrrr@{}}
\toprule
order & \(a\) & complex \(\sinh\) & Fourier & SIREN & MscaleDNN \\
\midrule
4 & 2 & 4.45e-1 & \textbf{3.65e-1} & 5.55e-1 & 1.30e0 \\
4 & 3 & \textbf{4.10e-1} & 5.90e-1 & 7.68e-1 & 1.51e0 \\
6 & 2 & \textbf{5.60e-1} & 8.02e-1 & 9.74e-1 & 1.33e0 \\
6 & 3 & \textbf{8.34e-1} & 9.11e-1 & 9.99e-1 & 1.57e0 \\
\bottomrule
\end{tabular}
\end{table}

\begin{figure}[htbp]
  \centering
  \includegraphics[width=\linewidth]{fig_cahn}
  \caption{Cahn--Hilliard. (a) Relative \(L^2\) versus amplitude \(a\) at sixth
  order; (b)--(c) relative \(L^2\) and loss versus training time at the
  sixth-order, \(a=2\) case. Bands show the seed min--max.}
  \label{fig:cahn}
\end{figure}

The nonlinear high-order phase field returns to the favourable regime. The
complex network wins three of the four instances, including both sixth-order
cases, and trails a Fourier network only on the smoothest fourth-order amplitude
(\(a=2\)). MscaleDNN diverges on all four. The result shows the advantage on
balanced high-order operators survives the addition of a cubic nonlinearity,
which the Taylor-jet backend assembles exactly from \eqref{eq:lap_u3}.
\FloatBarrier

% ===========================================================================
\subsection{Cubic nonlinear Schr\"odinger}\label{sec:bm:nls}

The first genuinely complex-valued benchmark is the cubic nonlinear
Schr\"odinger (NLS) equation,
\begin{equation}\label{eq:nls_pde}
  \mathrm{i}\,u_t + \tfrac12\,u_{xx} + |u|^2 u = f,
  \qquad u:\ \Omega \to \mathbb{C},
\end{equation}
with the standing bright-soliton family as the manufactured solution,
\(u_\star(x,t) = \operatorname{sech}(x)\,\mathrm{e}^{\mathrm{i}kt}\). A direct
computation gives
\begin{align}\label{eq:nls_source}
  \mathrm{i}u_{\star,t} + \tfrac12 u_{\star,xx} + |u_\star|^2 u_\star
  &= \big(\tfrac12 - k\big)\,u_\star,
\end{align}
so \(f=(\tfrac12-k)u_\star\), and at \(k=\tfrac12\) the soliton is exact. The
temporal frequency is swept over \(k\in\{1,2,4\}\). A single complex-sinh
network carries \(u\) in \(\mathbb{C}\) natively; the real baselines carry a
split real/imaginary pair. This is the complex-valued benchmark of Raissi et
al.~\cite{Raissi2019}, whose canonical setting is
\(\mathrm{i}h_t + \tfrac12 h_{xx} + |h|^2 h = 0\) on \(x\in[-5,5]\),
\(t\in[0,\pi/2]\); we use the same domain window with the manufactured standing
soliton so the reference is known in closed form.
\emph{Setting:} \(x\in[-5,5]\), \(t\in[0,\pi/2]\) (inputs normalized to
\([-1,1]^2\)); Dirichlet/initial penalty; complex vs.\ split-real.

\begin{table}[htbp]
\centering
\small
\caption{Cubic NLS (complex-valued): relative \(L^2\) error \eqref{eq:relL2}.
Real baselines carry a split real--imaginary pair.}
\label{tab:res-nls}
\begin{tabular}{@{}lrrrr@{}}
\toprule
\(k\) & complex \(\sinh\) & SIREN & Fourier & split \(\tanh\) \\
\midrule
1 & \textbf{7.56e-3} & 1.20e-2 & 1.82e-2 & 2.76e-1 \\
2 & \textbf{3.26e-2} & 7.06e-2 & 7.90e-2 & 3.84e-1 \\
4 & \textbf{4.13e-2} & 8.82e-2 & 8.95e-2 & 3.95e-1 \\
\bottomrule
\end{tabular}
\end{table}

\begin{figure}[htbp]
  \centering
  \includegraphics[width=\linewidth]{fig_nls}
  \caption{Cubic NLS. (a) Relative \(L^2\) versus temporal frequency \(k\);
  (b)--(c) relative \(L^2\) and loss versus training time at \(k=2\). Bands show
  the seed min--max.}
  \label{fig:nls}
\end{figure}

The complex network is the most accurate method at every \(k\), leading the best
split-real network by \(1.6\)--\(2.2\times\) and reaching a relative error below
\(10^{-2}\) at \(k=1\). The split \(\tanh\) pair never falls below \(0.27\). This
is the regime where a complex value field is most natural -- a split
real--imaginary surrogate must coordinate two independent networks to represent a
single analytic field -- and the advantage is correspondingly the most
consistent in the suite.
\FloatBarrier

% ===========================================================================
\subsection{Time-harmonic Maxwell (lossy medium)}\label{sec:bm:maxwell}

The second complex-valued benchmark is the time-harmonic transverse-magnetic
(TM\(_z\)) Maxwell system, which for the out-of-plane field reduces to a complex
Helmholtz equation,
\begin{equation}\label{eq:maxwell_pde}
  \Delta E + \kappa^2 E = f \quad \text{in } (-1,1)^2,
  \qquad \kappa^2 = (a\pi)^2(1+\mathrm{i}\beta),
  \qquad \beta = 0.2,
\end{equation}
with a traveling-plane-wave solution \(E_\star = \mathrm{e}^{\mathrm{i}a\pi(x+y)}\),
Dirichlet data \(E=E_\star\), and source
\(f=(-2(a\pi)^2+\kappa^2)E_\star\), swept over \(a\in\{2,4,6\}\). The complex
permittivity makes the field genuinely complex-valued. Time-harmonic Maxwell /
frequency-domain scattering is an inherently complex-valued family for which
complex-valued surrogates show clear advantages over real
ones~\cite{Jiang2024Maxwell}; the lossy \(\mathrm{i}\beta\) term is the linear
complex companion to the cubic NLS test.
\emph{Setting:} \(\Omega=(-1,1)^2\); Dirichlet \(E=E_\star\); \(\beta=0.2\);
\(\omega_0=\max(10,2\pi a)\), \(\sigma=\max(2,\pi a)\); complex vs.\ split-real.

\begin{table}[htbp]
\centering
\small
\caption{Time-harmonic Maxwell, lossy medium (complex-valued): relative \(L^2\)
error \eqref{eq:relL2}.}
\label{tab:res-maxwell}
\begin{tabular}{@{}lrrrr@{}}
\toprule
\(a\) & complex \(\sinh\) & SIREN & Fourier & split \(\tanh\) \\
\midrule
2 & \textbf{2.31e-1} & 3.63e-1 & 3.94e-1 & 1.76e0 \\
4 & \textbf{4.84e-1} & 7.62e-1 & 8.29e-1 & 1.82e0 \\
6 & \textbf{6.22e-1} & 9.38e-1 & 1.02e0  & 1.55e0 \\
\bottomrule
\end{tabular}
\end{table}

\begin{figure}[htbp]
  \centering
  \includegraphics[width=\linewidth]{fig_maxwell}
  \caption{Time-harmonic Maxwell (lossy). (a) Relative \(L^2\) versus
  loss/frequency scale \(a\); (b)--(c) relative \(L^2\) and loss versus training
  time at \(a=4\). Bands show the seed min--max.}
  \label{fig:maxwell}
\end{figure}

As in the NLS case the complex network is the most accurate method on every
instance, by \(1.5\)--\(1.6\times\) over the best split-real network, while the
split \(\tanh\) pair fails outright (errors above \(1.5\)). The lossy term forces
a genuinely complex field through a linear second-order operator, isolating the
representational benefit of a native complex value field from any nonlinearity.
\FloatBarrier

% ===========================================================================
\subsection{Summary and acceptance criterion}\label{sec:osc-summary}

Table~\ref{tab:res-summary} aggregates the contest by family. Counting an
instance as a win when the complex network attains the strictly smallest
relative \(L^2\) error, the complex \(\sinh\) ansatz is the best of the field on
\(32\) of the \(45\) instances. The wins concentrate on the regimes the ansatz
is designed for -- balanced high-order operators (polyharmonic, plate,
Cahn--Hilliard) and genuinely complex-valued fields (NLS and Maxwell, won at
every instance) -- and on non-separable high-frequency targets (the chirp);
the gains shrink or reverse on the first-order-in-time dispersive KdV, on the
one-dimensional beam and low-frequency plate modes, and on the low-frequency or
strongly directional Helmholtz problems.

We had set the stronger \emph{acceptance criterion} that the advantage factor
\eqref{eq:adv} grow monotonically with order or wavenumber. The data do not
support this in its strict form: although the per-family advantage is almost
always above one, it does not increase monotonically, and at the most extreme
order or wavenumber every method saturates near
\(\mathcal{E}_{L^2}\approx 1\), so the ratio is no longer informative. The honest
reading is therefore the weaker but robust one: \emph{at matched parameters and
a shared protocol, a complex value field is competitive-to-superior across the
oscillatory suite and clearly superior on balanced high-order and complex-valued
problems, but it is not a universal win and does not overcome the budget-limited
saturation of high-frequency PINNs.}

\begin{table}[htbp]
\centering
\small
\caption{Per-family summary. ``wins'' counts instances on which the complex
\(\sinh\) network attains the smallest relative \(L^2\) error; ``adv'' is the
range of the advantage factor \eqref{eq:adv} over the family's instances.}
\label{tab:res-summary}
\begin{tabular}{@{}llrl@{}}
\toprule
family & regime & wins & adv range \\
\midrule
Helmholtz (high \(k\))       & high-freq.\ 2nd order   & 5/5 & 1.07--1.55 \\
Helmholtz (anis.\ \(1,4\))   & directional 2nd order   & 0/1 & 0.57 \\
Helmholtz (var.\ \(\kappa\)) & heterogeneous 2nd order & 1/3 & 0.59--1.13 \\
radial chirp                 & non-separable 2nd order & 4/4 & 1.15--1.88 \\
\midrule
polyharmonic 1D              & order sweep \(2m\le10\)  & 4/5 & 0.88--5.94 \\
polyharmonic 2D              & order sweep \(2m\le6\)   & 2/3 & 0.80--12.96 \\
plate (isotropic)            & 4th order, modes        & 2/3 & 0.92--1.85 \\
plate (mixed modes)          & 4th order, anisotropic  & 3/3 & 1.42--1.77 \\
beam (1D)                    & 4th order, 1D           & 1/3 & 0.31--1.22 \\
linearized KdV               & 3rd order, dispersive   & 1/5 & 0.72--1.07 \\
Cahn--Hilliard               & 4th/6th, nonlinear      & 3/4 & 0.82--1.44 \\
\midrule
cubic NLS                    & complex-valued          & 3/3 & 1.59--2.16 \\
Maxwell (lossy)              & complex-valued          & 3/3 & 1.51--1.57 \\
\midrule
\textbf{total}               &                         & \textbf{32/45} & \\
\bottomrule
\end{tabular}
\end{table}
\FloatBarrier
